\documentclass[10pt,a4paper]{article}
\usepackage[T1]{fontenc}
\usepackage[utf8]{inputenc}
\usepackage{lmodern}
\usepackage[english]{babel}
\usepackage[landscape,margin=15mm]{geometry}
\usepackage{booktabs,array,graphicx,amssymb}
\usepackage[hidelinks]{hyperref}

\setlength{\parindent}{0pt}
\setlength{\tabcolsep}{3.5pt}
\renewcommand{\arraystretch}{1.4}
% booktabs-Abstände auf 0, sonst haben die senkrechten Linien an jeder Regel eine Lücke.
\setlength{\aboverulesep}{0pt}
\setlength{\belowrulesep}{0pt}
\newcolumntype{C}{>{\centering\arraybackslash}m{4.5mm}}
\newcolumntype{P}[1]{>{\raggedright\arraybackslash}m{#1}}
\newcommand{\rot}[1]{\rotatebox{90}{#1}}
\newcommand{\y}{\checkmark}

\begin{document}

{\large Related work matrix}
\vspace{2ex}

\footnotesize
\begin{tabular}{@{}P{30mm}|CCCCCP{22mm}|CCCCCP{24mm}|CCCCCC|P{17mm}C@{}}
\toprule
& \multicolumn{6}{c|}{\textbf{Frameworks}} & \multicolumn{6}{c|}{\textbf{Use cases}} & \multicolumn{6}{c|}{\textbf{Metrics}} & & \\
\textbf{Paper} &
\rot{XDP} & \rot{AF\_XDP} & \rot{DPDK} & \rot{Snabb} & \rot{Linux stack} & Other &
\rot{Drop} & \rot{L2FWD} & \rot{L3FWD} & \rot{Encap/Decap} & \rot{SFC} & Other &
\rot{Throughput} & \rot{Latency} & \rot{Jitter} & \rot{CPU} & \rot{Memory/cache} & \rot{Power} &
\textbf{Link rate} & \rot{Same code} \\
\midrule
Gallenmüller et al.\ (2015)~\cite{gallenmueller2015} &
 & & \y & & & netmap, PF\_RING ZC &
 & \y & & & & emulated CPU load, random memory access &
\y & \y & & \y & \y & &
10G & \\
\midrule
Høiland-Jørgensen et al.\ (2018)~\cite{xdp2018} &
\y & & \y & & \y & IPVS &
\y & \y & \y & \y & & DoS filter, L4 load balancing &
\y & \y & & \y & & &
100G\textsuperscript{c} & \\
\midrule
Hohlfeld et al.\ (2019)~\cite{hohlfeld2019} &
\y & \y\textsuperscript{a} & & & & SmartNIC offload (XDP) &
\y & & & & & emulated CPU load, DNS responder, map lookup/update &
\y & \y & & \y & & &
10G & \y\textsuperscript{b} \\
\bottomrule
\end{tabular}

\vspace{1.5ex}
\scriptsize
% Als Tabelle, damit alle Zeilen trotz Hochstellung gleich hoch sind.
{\renewcommand{\arraystretch}{1.2}%
\begin{tabular}{@{}l@{~}l@{}}
\multicolumn{2}{@{}l}{\y~measured \quad empty = not covered} \\
\textsuperscript{a} & Generic (copy) mode only, the NIC driver had no driver-specific AF\_XDP mode on Linux 4.18. \\
\textsuperscript{b} & Same XDP program at all three execution points. \\
\textsuperscript{c} & At 64\,B limited by PCIe: 115\,Mpps drop, 70\,Mpps forwarding out the same NIC. \\
\end{tabular}}

\begin{thebibliography}{3}
\scriptsize
\bibitem{gallenmueller2015}
S.~Gallenmüller, P.~Emmerich, F.~Wohlfart, D.~Raumer and G.~Carle.
\emph{Comparison of Frameworks for High-Performance Packet IO}. ANCS, 2015.
\url{https://ieeexplore.ieee.org/document/7110118}

\bibitem{xdp2018}
T.~Høiland-Jørgensen, J.~D.~Brouer, D.~Borkmann, J.~Fastabend, T.~Herbert, D.~Ahern and D.~Miller.
\emph{The eXpress Data Path: Fast Programmable Packet Processing in the Operating System Kernel}. CoNEXT, 2018.
\url{https://dl.acm.org/doi/10.1145/3281411.3281443}

\bibitem{hohlfeld2019}
O.~Hohlfeld, J.~Krude, J.~H.~Reelfs, J.~Rüth and K.~Wehrle.
\emph{Demystifying the Performance of XDP BPF}. NetSoft, 2019.
\url{https://ieeexplore.ieee.org/document/8806651}
\end{thebibliography}

\end{document}
